\par\SourceLeaf{421}%
\noindent\hspace*{1.2em}Beim Fortschreiten entlang der Kurve bleiben also die Seite, zu der sie konvex ist, oder die Vorzeichen in den obigen Formeln unverändert, solange
\[
\frac{\delta V}{\delta\varrho}
\]
sein Vorzeichen nicht ändert. Sofern $V$ eine stetige Funktion ist, kann ein solcher Vorzeichenwechsel erst nach einem Durchgang durch den Wert $0$ eintreten. Das setzt aber notwendig voraus, dass zugleich $p$ und $q=0$ werden. An einer solchen Stelle wird der Krümmungsradius unendlich groß, oder die Krümmung verschwindet. Da dann hier, im Allgemeinen gesprochen,
\[
-p\sin\varphi+q\cos\varphi
\]
sein Vorzeichen ändert, liegt hier ein Wendepunkt vor.

\Needspace{5\baselineskip}
\begin{center}5.\end{center}
\noindent\hspace*{1.2em}Der Fall, in dem die Natur der Kurve dadurch ausgedrückt wird, dass $y$ einer gegebenen Funktion von $x$ gleichgesetzt wird, nämlich $y=X$, ist im vorigen Fall enthalten, wenn man
\[
V=X-y
\]
setzt. Setzt man
\[
dX=X'dx,\qquad dX'=X''dx,
\]
so erhält man
\[
\begin{gathered}
p=X',\qquad q=-1,\\
P=X'',\qquad Q=0,\qquad R=0,
\end{gathered}
\]
also
\[
\pm\frac1r=\frac{X''}{(1+X'X')^{\frac32}}.
\]
Da $q$ hier negativ ist, gilt für zunehmendes $x$ das obere Vorzeichen. Man kann es hier also kurz so ausdrücken: Bei positivem $X''$ ist die Kurve zu der Seite konkav, auf der die $y$-Achse gegenüber der $x$-Achse liegt. Bei negativem $X''$ ist die Kurve zu dieser Seite konvex.
